%---------- Inleiding ---------------------------------------------------------

% TODO: Is dit voorstel gebaseerd op een paper van Research Methods die je
% vorig jaar hebt ingediend? Heb je daarbij eventueel samengewerkt met een
% andere student?
% Zo ja, haal dan de tekst hieronder uit commentaar en pas aan.

%\paragraph{Opmerking}

% Dit voorstel is gebaseerd op het onderzoeksvoorstel dat werd geschreven in het
% kader van het vak Research Methods dat ik (vorig/dit) academiejaar heb
% uitgewerkt (met medesturent VOORNAAM NAAM als mede-auteur).
%

\section{Inleiding}%
\label{sec:inleiding}

% Waarover zal je bachelorproef gaan? Introduceer het thema en zorg dat volgende zaken zeker duidelijk aanwezig zijn:



% \autocite{LiterateProgramming}

% \begin{itemize}
%   \item kaderen thema
%   \item de doelgroep
%   \item de probleemstelling en (centrale) onderzoeksvraag
%   \item de onderzoeksdoelstelling
% \end{itemize}

% Denk er aan: een typische bachelorproef is \textit{toegepast onderzoek}, wat betekent dat je start vanuit een concrete probleemsituatie in bedrijfscontext, een \textbf{casus}. Het is belangrijk om je onderwerp goed af te bakenen: je gaat voor die \textit{ene specifieke probleemsituatie} op zoek naar een goede oplossing, op basis van de huidige kennis in het vakgebied.

% De doelgroep moet ook concreet en duidelijk zijn, dus geen algemene of vaag gedefinieerde groepen zoals \emph{bedrijven}, \emph{developers}, \emph{Vlamingen}, enz. Je richt je in elk geval op it-professionals, een bachelorproef is geen populariserende tekst. Eén specifiek bedrijf (die te maken hebben met een concrete probleemsituatie) is dus beter dan \emph{bedrijven} in het algemeen.

% Formuleer duidelijk de onderzoeksvraag! De begeleiders lezen nog steeds te veel voorstellen waarin we geen onderzoeksvraag terugvinden.

% Schrijf ook iets over de doelstelling. Wat zie je als het concrete eindresultaat van je onderzoek, naast de uitgeschreven scriptie? Is het een proof-of-concept, een rapport met aanbevelingen, \ldots Met welk eindresultaat kan je je bachelorproef als een succes beschouwen?


% Wetenschappelijke instituten zoals ILVO (Instituut voor Landbouw-, Visserij- en voedingsOnderzoek) maken intensief gebruik van complexe rekenmodellen om emissies, productie-impact en andere landbouwparameters te analyseren. Deze modellen bestaan uit ingewikkelde verzamelingen van formules die door de onderzoekers worden ontwikkeld en vervolgens geimplementeerd worden door het IT departement in softwaretoepassingen.
% Aangezien de onderzoekers deze toepassingen gebruiken om hun hypotheses te testen en scenario's simuleren, is de kwaliteit en duidelijkheid van deze implementatie van groot belang. In praktijk blijken deze modellen moeilijk aan te passen zodra ze in de code zijn gegoten.

% Binnen ILVO wordt dit probleem duidelijk zichtbaar in hun .NET applicaties waar emissieformules en bedrijfsparameters zijn verwerkt in de C\# logica. Onderzoekers beschikken over de nodige inhoudelijke expertise om hun formules en parameters te wijzigen, maar hebben niet altijd genoeg programmeer kennis om deze wijzigingen zelfstandig door te voeren in de software. Voor elke wijziging zijn zij afhankelijk van het IT departement, waardoor kleine veranderingen vaak vertraging oplopen.
% Tegelijkertijd is de huidige documentatie niet altijd voldoende nauw verwoven met de code, waardoor het moeilijk is om te achterhalen waar specifieke berekeningen plaatsvinden en welke aannames aan de basis liggen.

% De doelgroep van dit onderzoek bestaat dan ook uit twee nauw samenwerkende groepen binnen ILVO:\@ de onderzoekers die de wetenschappelijke modellen ontwikkelen en de IT-professionals die deze modellen implementeren, onderhouden en documenteren. Beide groepen ervaren hinder door de huidige werkwijze: onderzoekers missen autonomie, terwijl IT'ers veel tijd spenderen aan kleine wijzigingen, uitleg en debugging.

% Daarom staat in deze bachelorproef de volgende centrale onderzoeksvraag centraal:
% Hoe kunnen we binnen een .NET-gebaseerde (C\# \& Razor Pages) onderzoeksapplicatie de kloof verkleinen tussen IT-technische logica en wetenschappelijke analyse, zodat onderzoekers van ILVO zonder diepgaande programmeerkennis inzicht krijgen in de werking van hun intern model en er zelf mee kunnen experimenteren?
% Deze vraag is relevant omdat ze zowel de transparantie van wetenschappelijke modellen verbetert als de samenwerking tussen onderzoek en IT efficienter maakt.

% Het doel van dit onderzoek is om een proof-of-concept te ontwikkelen waarin de wetenschappelijke onderzoekers met hun modellen kunnen experimenteren zonder dat zij een extensieve kennis van programmeren nodig hebben. Hierbij wordt onderzocht welke documentatie vormen het meest geschikt zijn om de interne werking van hun modellen begrijpelijk te maken. Daarnaast wordt nagegaan hoe parameters veilig kunnen worden aangepast. Het uiteindelijke resultaat moet een concreet voorstel en prototype opleveren dat ILVO kan gebruiken om hun onderzoeks modellen toegankelijker, transparanter en beter onderhoudbaar te maken.

% Dit onderzoek wordt bewust afgebakend: het richt zich niet op het herontwerpen van het volledige ILVO-model of het bouwen van een productieklare applicatie, maar wel op het ontwikkelen en evalueren van technieken, workflows en tooling die de interactie tussen code en onderzoek verbeteren. Theoretisch kader wordt gezocht binnen software engineering, documentatieonderzoek, literate programming, domain-driven design en transparantie van wetenschappelijke modellen.


% \section{Inleiding}
% \label{sec:inleiding}

Wetenschappelijke instituten zoals ILVO (Instituut voor Landbouw-, Visserij- en Voedingsonderzoek) maken intensief gebruik van complexe rekenmodellen om emissies, productie-impact en andere landbouwparameters te analyseren. Deze modellen bestaan uit clusters van formules die door onderzoekers worden ontwikkeld en vervolgens door de IT-afdeling worden geïmplementeerd in softwaretoepassingen. Aangezien onderzoekers deze toepassingen gebruiken om hypothesen te testen en scenario’s te simuleren, is de kwaliteit, transparantie en aanpasbaarheid van deze implementaties van groot belang. In de praktijk blijken dergelijke modellen echter moeilijk wijzigbaar en weinig inzichtelijk zodra ze in code zijn gegoten.

Binnen ILVO wordt dit probleem concreet zichtbaar in .NET-gebaseerde applicaties waarin emissieformules en bedrijfsparameters zijn verwerkt in C\#-logica. Onderzoekers beschikken over diepgaande inhoudelijke expertise om hun modellen en aannames te verfijnen, maar hebben niet altijd voldoende programmeerkennis om deze wijzigingen zelfstandig in de software door te voeren. Voor elke aanpassing zijn zij afhankelijk van de IT-afdeling, waardoor zelfs beperkte experimenten of kleine parameterwijzigingen vertraging oplopen. Bovendien is de bestaande documentatie niet altijd nauw verweven met de code, wat het moeilijk maakt om te achterhalen waar specifieke berekeningen plaatsvinden en welke aannames eraan ten grondslag liggen.

De doelgroep van dit onderzoek bestaat uit twee nauw samenwerkende groepen binnen ILVO: enerzijds de onderzoekers die de wetenschappelijke modellen ontwikkelen, en anderzijds de IT-professionals die deze modellen implementeren, onderhouden en documenteren. Beide groepen ondervinden hinder van de huidige werkwijze: onderzoekers ervaren een gebrek aan autonomie en inzicht, terwijl IT-professionals disproportioneel veel tijd besteden aan kleine aanpassingen, ondersteuning en foutopsporing.

Vanuit deze concrete probleemsituatie wordt in deze bachelorproef de volgende centrale onderzoeksvraag geformuleerd:
\textit{Hoe kunnen we binnen een .NET-gebaseerde (C\# \& Razor Pages) onderzoeksapplicatie de kloof verkleinen tussen IT-technische logica en wetenschappelijke analyse, zodat onderzoekers van ILVO zonder diepgaande programmeerkennis inzicht krijgen in de werking van hun intern model en er zelf mee kunnen experimenteren?}

Om deze onderzoeksvraag te kunnen beantwoorden, wordt het onderzoek opgesplitst in een aantal samenhangende deelvragen die zowel het probleemdomein als het oplossingsdomein bestrijken. Eerst wordt onderzocht hoe de huidige samenwerking en workflow tussen onderzoekers en IT-professionals verloopt en waar deze in de praktijk vastloopt. Daarnaast wordt nagegaan welke onderdelen van het bestaande .NET-model en de bijhorende code moeilijk te begrijpen of aan te passen zijn voor onderzoekers. Ook wordt onderzocht in welke mate de huidige documentatie tekortschiet in het zichtbaar maken van berekeningen, aannames en afhankelijkheden binnen het model.

Op basis van deze probleemanalyse richt het onderzoek zich vervolgens op mogelijke oplossingsrichtingen. Daarbij wordt onderzocht welke documentatie- en interactievormen onderzoekers ondersteunen bij het begrijpen van wetenschappelijke modellen, zoals inline annotaties, interactieve documentatie of computationele notebooks. Verder wordt nagegaan hoe onderzoekers veilig kunnen experimenteren met modelparameters en formules, bijvoorbeeld via simulaties of gevoeligheidsanalyses, zonder de onderliggende softwarestructuur te doorbreken. Tot slot wordt onderzocht hoe documentatie en code nauwer met elkaar verweven kunnen worden zodat wijzigingen aan het model automatisch transparant blijven voor alle betrokkenen.

Het doel van dit onderzoek is het ontwikkelen van een proof-of-concept waarin documentatie, modelberekeningen en experimenteerruimte voor onderzoekers dichter bij elkaar worden gebracht. Dit proof-of-concept moet aantonen hoe onderzoekers op een gecontroleerde en begrijpelijke manier met hun modellen kunnen werken, terwijl de onderhoudbaarheid en technische kwaliteit van de software behouden blijft. Het uiteindelijke resultaat bestaat uit een concreet voorstel en prototype dat ILVO kan ondersteunen bij het toegankelijker, transparanter en duurzamer maken van hun onderzoeksmodellen.

Dit onderzoek wordt bewust afgebakend: het richt zich niet op het herontwerpen van het volledige ILVO-model of het bouwen van een productieklare applicatie, maar op het ontwikkelen en evalueren van technieken, workflows en tooling die de interactie tussen code en onderzoek verbeteren. Het theoretisch kader wordt gezocht binnen software engineering, documentatieonderzoek, literate programming, domain-driven design en transparantie van wetenschappelijke modellen.


%---------- Stand van zaken ---------------------------------------------------

\section{Literatuurstudie}%
\label{sec:literatuurstudie}

%% Hier beschrijf je de \emph{state-of-the-art} rondom je gekozen onderzoeksdomein, d.w.z.\ een inleidende, doorlopende tekst over het onderzoeksdomein van je bachelorproef. Je steunt daarbij heel sterk op de professionele \emph{vakliteratuur}, en niet zozeer op populariserende teksten voor een breed publiek. Wat is de huidige stand van zaken in dit domein, en wat zijn nog eventuele open vragen (die misschien de aanleiding waren tot je onderzoeksvraag!)?

%% Je mag de titel van deze sectie ook aanpassen (literatuurstudie, stand van zaken, enz.). Zijn er al gelijkaardige onderzoeken gevoerd? Wat concluderen ze? Wat is het verschil met jouw onderzoek?

%% %% Verwijs bij elke introductie van een term of bewering over het domein naar de vakliteratuur, bijvoorbeeld~\autocite{Hykes2013}! Denk zeker goed na welke werken je refereert en waarom.
%% Draag zorg voor correcte literatuurverwijzingen! Een bronvermelding hoort thuis \emph{binnen} de zin waar je je op die bron baseert, dus niet er buiten! Maak meteen een verwijzing als je gebruik maakt van een bron. Doe dit dus \emph{niet} aan het einde van een lange paragraaf. Baseer nooit teveel aansluitende tekst op eenzelfde bron.

%% Als je informatie over bronnen verzamelt in JabRef, zorg er dan voor dat alle nodige info aanwezig is om de bron terug te vinden (zoals uitvoerig besproken in de lessen Research Methods).

%% % Voor literatuurverwijzingen zijn er twee belangrijke commando's:
%% % \autocite{KEY} => (Auteur, jaartal) Gebruik dit als de naam van de auteur
%% %   geen onderdeel is van de zin.
%% % \textcite{KEY} => Auteur (jaartal)  Gebruik dit als de auteursnaam wel een
%% %   functie heeft in de zin (bv. ``Uit onderzoek door Doll & Hill (1954) bleek
%% %   ...'')

%% Je mag deze sectie nog verder onderverdelen in subsecties als dit de structuur van de tekst kan verduidelijken.
\subsection{Literate programming}
Literate programming \autocite{LiterateProgramming} werd geïntroduceerd door \textcite{Knuth1984} als een methode waarbij code en documentatie 1 geïntegreerd geheel vormen, eerder dan gescheiden artefacten. De bedoeling is dat er 1 bestand is met de code en de documentatie errond. De tools halen de code eruit en creeeren een uitvoerbaar bestand en een ander bestand met de documentatie, dat kan een html web pagina zijn, of een \LaTeX \@ bestand / pdf. Het kan ook een Markdown bestand zijn.

In plaats van `code with comments' vertrekt literate programming vanuit het idee dat een programma in de eerste plaats een verhaal is dat aan een menselijke lezer moet worden uitgelegd, waarbij uitvoerbare code een secundaire afgeleide vorm is \autocite{Knuth1984}.
Sindsdien zijn er talrijke systemen ontstaan die Knuths principes toepassen in een moderne ontwikkelingsomgeving.

Hoewel Knuths oorspronkelijke WEB en CWEB nog weinig gebruikt worden in de hedendaagse softwarepraktijk, leven zijn ideeen voort in computationele notebooks, literate programminglibraries en research pipelines. Tools zoals Jupyter Notebooks, Org-mode Babel en Quarto implementeren een moderne interpretatie van Knuths `weave' en `tangle'-proces, waarin broncode zowel als uitvoerbaar script en als leesbare documentatie dient.

De execution haalt de source code eruit en zorgt voor een uitvoerbaar bestand. De Weave operatie zorgt er voor dat de documentatie uit het source bestand wordt gehaald en die in het gewenste formaat zet. De Tangle operatie verspreidt de source code in de respectievelijke bestanden zelf, niet direct uitvoerbaar.

\subsection{Org-mode Babel}
Org mode is een `major mode' voor GNU Emacs. Org mode kan voor heel veel gebruikt worden: van notities nemen tot todo-lists, computational notebooks en literate programming \autocite{orgmode2025}.

In Org Mode met Babel zijn er `src-blocks`:

\begin{minted}{text}
* Src-blocks in org-mode
#+BEGIN_SRC python :results verbatim
  nummer: int = 10
  return nummer
#+END_SRC

#+RESULTS:
: 10
\end{minted}


Hierin kan code worden uitgevoerd terwijl er tekst rond staat.
Dit is het idee van literate programming.
De meeste `cellen' in org-mode hebben hun eigen omgeving, maar er bestaan manieren om dit te vermijden.
Babel verbetert de src-blocks door een paar zaken te voorzien:
\begin{itemize}
	\item interactieve en programma executie van code blokken
	\item code blokken die functies met parameters hebben, kunnen andere code blokken accepteren en oproepen, en
	\item bestanden exporteren voor literate programming
\end{itemize}

\subsection{Jupyter notebook}
Jupyter notebooks zijn gemaakt voor python. Deze notebooks worden door vele mensen gebruikt en in verschillende omgevingen.
Dit kan gaan van een student die python leert en zijn notities er bij wil schrijven tot een volledige data analyse door professionele onderzoekers.
Dit is een standaard voor computationeel onderzoek en datascience \autocite{jupyternotebook}.
Het werkt met verschillende `cellen`: de markdown cellen en de python cellen.
De python cellen verbinden met een virtuele omgeving en voeren hierop telkens hun code uit.
Aangezien ze werken op dezelfde omgeving zijn de cellen verbonden aan elkaar.
Het grootste verschil met org-mode is dat een \text{.}org bestand een plain tekst bestand is. Het is mogelijk om dit bestand te bekijken in eender welke editor.
Jupyter notebooks niet. Je hebt een editor nodig die jupyter notebooks ondersteunt. Het bekentste is VsCode van Microsoft.
Er bestaat ook de officiele command en app van jupyter zelf die een server opstart en dan in de browser beschikbaar is.

Maar jupyter notebooks hebben zeker hun minpunten.
\textcite{Rule2018} tonen aan dat Jupyter Notebooks niet-lineaire uitvoergeschiedenis hebben, wat reproduceerbaarheid ernstig bemoeilijkt.
Ook omdat de cellen in een willekeurige volgorde kunnen worden uitgevoerd, moet er opgelet worden met de run volgorde.
De notebooks hebben ook een onzichtbare global state.
Het zijn json bestanden, dus versie controle beheer kan hier ook moeilijk mee doen.
Aangezien dat alles ook in 1 bestand zit, kan het snel `bloated' zijn.

\subsection{RMarkdown}
\texttt{RMarkdown} \autocite{RStudio2020} is ook een computational notebook zoals Jupyter Notebooks,
maar het ondersteunt meerdere talen. Het volgt een structuur die sterk aansluit bij de principes van
literate programming. Een \texttt{.Rmd}-bestand wordt verwerkt door \texttt{knitr} \autocite{Xie2025},
een transparante engine voor het genereren van dynamische rapporten in R, ontwikkeld binnen de
softwareomgeving voor statistische computing \texttt{R} \autocite{RFoundation}.

\texttt{Knitr} voert alle code blokken uit en maakt een nieuw Markdown bestand aan met alle code en hun uitkomsten.
Het gemaakte Markdown bestand wordt dan geevalueerd door pandoc \autocite{pandoc}, die dan het finale product maakt.
Vb:
\begin{minted}{text}
---
title: "viridis demo"
output: html_document
---

```{r include - FALSE}
library(viridis)
```
The code below ...

## Colors
```{r}
image(volcano, col - viridis(200))
```
\end{minted}

\subsection{Tussenbesluit}



%---------- Methodologie ------------------------------------------------------
\section{Methodologie}%
\label{sec:methodologie}

% Hier beschrijf je hoe je van plan bent het onderzoek te voeren. Welke onderzoekstechniek ga je toepassen om elk van je onderzoeksvragen te beantwoorden? Gebruik je hiervoor literatuurstudie, interviews met belanghebbenden (bv.~voor requirements-analyse), experimenten, simulaties, vergelijkende studie, risico-analyse, PoC, \ldots?

% Valt je onderwerp onder één van de typische soorten bachelorproeven die besproken zijn in de lessen Research Methods (bv.\ vergelijkende studie of risico-analyse)? Zorg er dan ook voor dat we duidelijk de verschillende stappen terug vinden die we verwachten in dit soort onderzoek!

% Vermijd onderzoekstechnieken die geen objectieve, meetbare resultaten kunnen opleveren. Enquêtes, bijvoorbeeld, zijn voor een bachelorproef informatica meestal \textbf{niet geschikt}. De antwoorden zijn eerder meningen dan feiten en in de praktijk blijkt het ook bijzonder moeilijk om voldoende respondenten te vinden. Studenten die een enquête willen voeren, hebben meestal ook geen goede definitie van de populatie, waardoor ook niet kan aangetoond worden dat eventuele resultaten representatief zijn.

% Uit dit onderdeel moet duidelijk naar voor komen dat je bachelorproef ook technisch voldoen\-de diepgang zal bevatten. Het zou niet kloppen als een bachelorproef informatica ook door bv.\ een student marketing zou kunnen uitgevoerd worden.

% Je beschrijft ook al welke tools (hardware, software, diensten, \ldots) je denkt hiervoor te gebruiken of te ontwikkelen.

% Probeer ook een tijdschatting te maken. Hoe lang zal je met elke fase van je onderzoek bezig zijn en wat zijn de concrete \emph{deliverables} in elke fase?


% \subsection{Probleem- en domeinanalyse}
% In deze fase wordt onderzocht hoe ILVO momenteel werkt met hun \text{.NET}-gebaseerde modellen.


% \begin{itemize}
% \item interviews en overlegmomenten met de onderzoekers (inhoudelijke experts) en de ontwikkelaars
% \item Analyse van bestaande \texttt{C\#}-code, Razor pages, berekeningslogica en interne documentatie
% \item Identificeren van
%   \begin{itemize}
%   \item waar berekeningen plaatsvinden
%   \item welke parameters en formules onderzoekers willen aanpassen
%   \item typische fouten of bottlenecks
%   \item waarvoor onderzoekers nu afhankelijk zijn van IT.\@
%   \end{itemize}
% \end{itemize}

% \subsection{Literatuurstudie \& technologieverkenning}
% Deze fase biedt het theoretische en technische kader dat rest van het onderzoek ondersteunt.

% \begin{itemize}
% \item Literate programming (\autocite{Knuth1984}, morderne interpretaties in software engineering)
% \item Reproducible research workflows in wetenschappelijke software
% \item Documentation-as-code: DocFX, Jupyter notebook, quarto, org mode
% \item Computational notebooks: jupyter, \text{.NET} Interactive, polyglot notebooks
% \item annotatie- en traceertools binnen codebases
% \item DSL's en transparantie in wetenschappelijke modellen
% \end{itemize}

% \subsection{Proof-of-Concept ontwikkeling}
% Op basis van de analyse en literatuur wordt minstens 1, maar mogelijk meerdere POC's ontwikkeld.
% Het doel is niet een productieklare applicatie, maar het aantonen van mogelijke oplossingsrichtingen.

% Mogelijke richtingen voor POC's zijn:

% \begin{enumerate}
% \item Interactieve notebooks:
%   De formules zijn snel aan te passen en het resultaat is direct te zien. Er is een automatische uitleg van parameters. Er zijn ook visualisaties van impact en gevoeligheidanalyses
% \item Inline annotatie-omgeving van C\#-code.
%   De commentaar blokken en uitleg wordt automatisch weergegeven naast de logicablokken. De documentatie gegenereerd via DocFX wordt teruggekoppeld naar de onderzoekers.
% \item Model viewer binnen een razor page.
%   Zorgt voor een grafisch overzicht van het model.
%   Mogelijks parameters met sliders.
%   Bedoelt om toekomstige scenario's te simuleren
% \item Koppeling tussen code en documentatie.
%   Alles in 1 bestand.
%   Eventueel via Quarto of org-mode.

% \end{enumerate}

% \subsection{Evaluatie en validatie}
% Er zullen korte gebruikerstesten plaatsvinden:
% \begin{itemize}
% \item Kunnen de onderzoekers formules terugvinden?
% \item Begrijpen ze de werking van het model beter?
% \item Kunnen ze parameters aanpassen zonder code te breken?
% \end{itemize}

% Semi-gestructureerde interviews met:
% \begin{itemize}
% \item Onderzoekers over bruikbaarheid
% \item IT'ers over impact op onderhoudbaarheid
% \end{itemize}

% Analyse van feedback:
% \begin{itemize}
% \item wat werkt?
% \item wat werkt niet?
% \item welke uitbreidingen zijn nodig?
Binnen een concrete casus bij het Instituut voor Landbouw-, Visserij- en Voedingsonderzoek (ILVO) wordt een toegepast, onderwerpgericht onderzoeksopzet gevolgd. De focus ligt op het analyseren van de huidige samenwerking tussen onderzoekers en softwareontwikkelaars rond een \text{.NET}-gebaseerd rekenmodel, en op het uitwerken van een haalbare oplossingsrichting die deze samenwerking ondersteunt. De methodologie is opgebouwd uit vijf opeenvolgende fasen, waarin zowel het probleemdomein als het oplossingsdomein systematisch worden onderzocht.

\subsection{Eerste fase}
In de eerste fase wordt het probleemdomein in kaart gebracht. Hierbij wordt onderzocht hoe onderzoekers en ontwikkelaars momenteel samenwerken rond het bestaande model, waar verantwoordelijkheden liggen en op welke punten de workflows elkaar belemmeren. Deze analyse vertrekt vanuit concrete use-cases, zoals het aanpassen van emissieformules, het testen van alternatieve parameters en het uitvoeren van gevoeligheidsanalyses. De gegevens worden verzameld via documentanalyse en overlegmomenten met betrokkenen, met als doel inzicht krijgen in waar de kloof tussen wetenschappelijke logica en technische implementatie zich manifesteert.

\subsection{Tweede fase}
De tweede fase richt zich op het identificeren van noden en vereisten van beide doelgroepen. Voor onderzoekers gaat het onder meer om transparantie van berekeningen, experimenteerruimte en begrijpelijke documentatie; voor ontwikkelaars om onderhoudbaarheid, stabiliteit en controle over data en interfaces. Deze noden worden gestructureerd geformuleerd als functionele en niet-functionele requirements, die richtinggevend zijn voor het verdere onderzoek. Hierbij wordt expliciet rekening gehouden met organisatorische aspecten, zoals taakverdeling en communicatiestromen, in lijn met inzichten uit onder meer Conway's Law \autocite{conwayslaw}.

\subsection{Derde fase}
In de derde fase wordt het oplossingsdomein onderzocht aan de hand van een gerichte literatuurstudie en analyse van bestaande tools en workflows. Concepten zoals \textcite{LiterateProgramming}, computational notebooks ( bijvoorbeeld jupyter notebook \autocite{jupyternotebook} en quarto \autocite{quarto2025}), documentatieplatformen (zoals docFX \autocite{docfx}) en model life-cycle tools (zoals mlflow \autocite{mlflow}) worden bestudeerd op hun relevantie voor gelijkaardige samenwerkingsproblemen. Daarbij wordt nagegaan in welke mate deze benaderingen geschikt zijn binnen een \text{.NET}-context, en welke principes overdraagbaar zijn, ook wanneer de tools zelf technologisch niet rechtstreeks inzetbaar zijn.

\subsection{Vierde fase}
Op basis van de voorgaande fasen volgt in de vierde fase een selectie en ontwerp van een oplossingsrichting. De eerder geindentificeerde requirements worden gebruikt om mogelijke oplossingen af te wegen en te beperken tot een haalbare subset. Hierbij wordt bewust gekozen voor een beperkte scope, met focus op het verbeteren van transparantie, documentatie en experimenteermogelijkheden, zonder het bestaande model volledig te herontwerpen. Het resultaat van deze fase is een concreet concept voor een ondersteunende workflow of tool, afgestemd op de ILVO-casus.

\subsection{Vijfde fase}
In de vijfde en laatste fase wordt deze oplossingsrichting geconcretiseerd in een proof-of-concept. Dit prototype demonstreert hoe onderzoekers inzicht kunnen krijgen in modelberekeningen en parameters, en hoe zij op een gecontroleerde manier kunnen experimenteren zonder diepgaande programmeerkennis. Het proof-of-concept wordt geevalueerd aan de hand van de eerder gedefinieerde use-cases en requirements. De bevindingen uit deze evaluatie vormen de basis voor conclusies en aanbevelingen over toepasbaarheid en meerwaarde van de voorgestelde aanpak binnen ILVO en gelijkaardige onderzoekscontexten.
% \end{itemize}
%---------- Verwachte resultaten ----------------------------------------------
\section{Verwacht resultaat, conclusie}%
\label{sec:verwachte_resultaten}

%% Hier beschrijf je welke resultaten je verwacht. Als je metingen en simulaties uitvoert, kan je hier al mock-ups maken van de grafieken samen met de verwachte conclusies. Benoem zeker al je assen en de onderdelen van de grafiek die je gaat gebruiken. Dit zorgt ervoor dat je concreet weet welk soort data je moet verzamelen en hoe je die moet meten.

%% Wat heeft de doelgroep van je onderzoek aan het resultaat? Op welke manier zorgt jouw bachelorproef voor een meerwaarde?

%% Hier beschrijf je wat je verwacht uit je onderzoek, met de motivatie waarom. Het is \textbf{niet} erg indien uit je onderzoek andere resultaten en conclusies vloeien dan dat je hier beschrijft: het is dan juist interessant om te onderzoeken waarom jouw hypothesen niet overeenkomen met de resultaten.

\subsection{Probleem- en domeinanalyse}
Een helder overzicht van de \textit{current flow} (onderzoekers \rightarrow  IT \rightarrow  applicatie)

Een lijst van problemen, zoals:
\begin{itemize}
	\item beperkte transparantie
	\item onvoldoende verweven documentatie
	\item moeilijk traceerbaarheid van berekeningen
	\item lange feedbackloops
\end{itemize}

\subsection{Literatuurstudie \& technologieverkenning}
Een lijst van mogelijke technieken en tools om de kloof tussen code en onderzoek te verkleinen.
Ook bepalen welke technieken \textit{haalbaar en relevant} zijn voor ILVO.\@


\subsection{Proof-of-Concept ontwikkeling}
Een geintegreerde POC die aantoont hoe onderzoekers:
\begin{itemize}
	\item parameters kunnen wijzigen,
	\item resultaten kunnen testen,
\end{itemize}
zonder afhankelijk te zijn van IT voor kleine experimenten.

\subsection{Evaluatie en validatie}
Conclusie over welke aanpak het meest effectief is om de kloof te verkleinen en een concreet advies voor ILVO met mogelijke roadmap.
